% From The tocbibind package [author: Peter Wilson, Herries Press]:
% The tocbibind package can be used to add document elements like a 
% bibliography or an index to the Table of Contents. The package is designed 
% to work with the four standard book, report, article and proc classes, and to 
% a limited extent with the ltxdoc class.
% 
\usepackage{tocbibind}

%
% From [Wikipedia]:
% The package hyperref provides LaTeX the ability to create hyperlinks within 
% the document. It works with pdflatex and also with standard "latex" used with 
% dvips and ghostscript or dvipdfm to build a PDF file. For example: 
%	- \hyperref[mainlemma]{lemma \ref{mainlemma}}
%	- \url{http://www.wikibooks.org}
%	- \href{http://www.wikibooks.org}{wikibooks home}
%	- \href{mailto:my_address@wikibooks.org}{
%		\nolinkurl{my_address@wikibooks.org}}
% 
\usepackage{hyperref}
  \hypersetup{
  	colorlinks=false,
  	pdfborder=0 0 0,
  	linkcolor=blue,
  	citecolor=black,
  	bookmarksopen=false,
  	bookmarksnumbered=true,
  	pdfstartview=FitH,
  	pdfview=FitH
	}

\usepackage{fancyhdr}
\usepackage{amsmath,amssymb,amsfonts}
\usepackage[pdftex]{graphicx}
\usepackage{textcomp}
\usepackage{xcolor}
\usepackage{booktabs}
\usepackage{multirow}
\usepackage{enumitem}
\usepackage{tikz}
\usetikzlibrary{backgrounds}
\usetikzlibrary{shapes,arrows.meta,positioning,fit,calc,backgrounds}

% Algorithm packages (fixed - no conflict)
\usepackage{algorithm}
\usepackage{algpseudocode}

\usepackage[numbers,sort]{natbib}

\newcommand{\var}[2]{\newcommand{#1}{#2}}